% NeurIPS 2026 typography; the six proposal sections follow ECSE 552 requirements.
% Limit: about two pages, under 1,000 words, plus one page for references.
% Red text is for the team: each [TO WRITE] marks a paragraph to write, and the
% red bullets below it are starting ideas. Replace both with your own prose.
% Literature notes: research_notes/action-free-robot-learning/
\documentclass{article}
\PassOptionsToPackage{numbers,sort&compress}{natbib}
\usepackage[preprint]{neurips_2026}
\usepackage[utf8]{inputenc}
\usepackage[T1]{fontenc}
\usepackage{amsmath}
\usepackage{amssymb}
\usepackage{microtype}
\usepackage{xcolor}
\usepackage[hidelinks]{hyperref}

\newcommand{\todo}[1]{\textcolor{red}{[TO WRITE: #1]}}
% Starting ideas, rendered in red so they are never submitted by accident.
\newenvironment{ideas}
  {\color{red}\small\begin{itemize}\setlength{\itemsep}{0pt}}
  {\end{itemize}}

\title{Learning to Manipulate a New Object\\from Action-Free Videos}
\author{
  Group 9 \\
  McGill University \\
  ECSE 552 --- Project proposal
}

\begin{document}
\maketitle

\section{Background}
% Rubric: A1--A2.
\paragraph{Application and problem.}
\todo{A1--A2, about 150 words. The application, the problem, and why a solution
is needed.}
\begin{ideas}
  \item Robot imitation learning needs demonstrations paired with motor commands,
    which are expensive to collect; videos without actions are much cheaper.
  \item Setting: action-labeled demonstrations with one object (cube) and only
    action-free videos with a new object (elongated block).
  \item Question: can these videos teach a policy to manipulate the new object?
  \item Labels are withheld in simulation on purpose, to emulate video-only data.
  \item Inverse-dynamics pseudo-labeling already exists (BCO, VPT), but assumes
    labeled and unlabeled data come from the same domain; here the object changes.
  \item Frame the contribution as a controlled study of adaptation across object
    geometry, not a new way of learning from video.
\end{ideas}
% Citable precedents (keys in ../references.bib):
%   torabi2018bco      -- BCO: IDM trained on known actions labels observation-only demos, then BC.
%   baker2022vpt       -- VPT: IDM on a small labeled set pseudo-labels large unlabeled video.
%   kim2023givinghand  -- IDM from robot play labels human eye-in-hand video (arXiv 2023).
%   collins2025amplify -- motion tokens from video + IDM; video-seen, action-unseen tasks (arXiv 2025).

% Rubric: A3.
\paragraph{Design constraints.}
\todo{A3, about 80 words. Requirements the design must satisfy, not team time or
equipment.}
\begin{ideas}
  \item Target data are only ordered RGB images: no actions, robot states, object
    poses, or annotations.
  \item The learned policy sees only current and past images.
  \item It must succeed in closed loop, with a fixed budget of source labels.
  \item Simulator state is never an input to a learned model.
\end{ideas}

\section{Aims/Goals}
% Rubric: A4.
\paragraph{Scope.}
\todo{A4, about 80 words. What is in and out of scope.}
\begin{ideas}
  \item OCBench visual single-block task in MuJoCo, UR5e arm with Robotiq gripper.
  \item Source: the existing cube. Target: a new elongated-block variant, which
    needs adapted expert grasps.
  \item Keep cameras, lighting, goal, and initial-position distribution fixed.
  \item Out of scope: human videos, real hardware, other spatial shifts.
\end{ideas}
% ocbench2026 -- OCBench (Park & Levine 2026, preprint): UR5e/Robotiq, MuJoCo, visual single-block lite task.

% Rubric: A5.
\paragraph{Hypothesis and contribution.}
\todo{A5, about 70 words. The hypothesis and the expected contribution.}
\begin{ideas}
  \item Domain-adversarial training of the IDM may give better target pseudo-labels,
    and therefore better policy success, than an unadapted IDM.
  \item Second question: does alignment help more in the IDM or in the policy?
  \item Risk to state honestly: invariance may also erase information needed to
    predict actions; a negative result is still informative.
\end{ideas}
% ganin2016dann      -- DANN: gradient reversal against a domain discriminator.
% zhao2019invariant  -- invariant features + low source error do not guarantee adaptation.

\section{Methodology}
% Rubric: A5--A6.
\paragraph{Data and implementation.}
\todo{A5--A6, about 200 words. Data generation, models, inputs and outputs,
training objectives.}
\begin{ideas}
  \item Generate demonstrations with OCBench's scripted expert and data generator;
    capture an image at every control step (the published visual data are sparse).
  \item Actions: OCBench's existing normalized joint and gripper increments.
  \item IDM: predicts the action from a short window of frames; it may look at
    future frames because it only labels data offline.
  \item Adaptation: a domain classifier (cube vs.\ block) with gradient reversal
    on the IDM encoder, while the action head keeps predicting source actions.
  \item Freeze the IDM, pseudo-label the block videos, train the policy on real
    cube labels plus block pseudo-labels.
  \item To decide: number of episodes, cameras, encoder, history length, loss
    weights.
\end{ideas}

% Rubric: A8.
\paragraph{Alternatives.}
\todo{A8, about 100 words. Alternatives run on our own data, and why ours could be
better.}
\begin{ideas}
  \item Source-only cloning (ignores the block videos).
  \item Pseudo-labels from an unadapted IDM (the BCO/VPT recipe).
  \item Alignment in the policy instead of the IDM.
  \item Reference trained with the true block actions (not a guaranteed upper bound).
  \item A published policy such as Diffusion Policy on the same data.
  \item Latent-action methods (LAPA, MotoVLA) as related work, too large to rerun.
\end{ideas}
% chi2023diffusionpolicy -- Diffusion Policy (RSS 2023); zhao2023act -- ACT (RSS 2023).
% ye2025lapa, spiridonov2025motovla -- latent-action / action-free pretraining (ICLR 2025, CoRL 2025).
% cheng2025gda -- GDA (NeurIPS 2025): OT alignment, uses action-labeled target demos.

% Rubric: A7.
\paragraph{Validation and evaluation.}
\todo{A7, about 100 words. Splits, model selection, metrics, episodes and seeds.}
\begin{ideas}
  \item Split by episode before cutting windows: source validation, target
    adaptation, target development, final test.
  \item Never use hidden block actions to select the adapted model.
  \item Select policies by rollout success, since validation loss need not track
    task success (OCBench).
  \item Metrics: success rate per seed, failure type (grasp, transport, placement),
    action error, source retention.
\end{ideas}
% ocbench2026 reports that validation loss need not track task success.

\section{Feasibility and resources}
% Rubric: A9.
\paragraph{Dataset, resources, and pilot.}
\todo{A9, about 80 words. Dataset, compute, pilot, timeline.}
\begin{ideas}
  \item Data are generated by us in OCBench on CPU MuJoCo; GPU access to confirm.
  \item Pilot checks: the expert succeeds on both objects, the source policy
    succeeds on the cube, success drops on the block, the reference learns the block.
\end{ideas}

\section{Risk analysis}
% Rubric: A10.
\paragraph{Pitfalls and mitigation.}
\todo{A10, about 80 words. What could fail and what we will do about it.}
\begin{ideas}
  \item The block is too hard for the expert: shorten it, narrow its rotations.
  \item Source-only already succeeds on the block: make the geometry change larger.
  \item The cube demos do not cover the motions the block needs.
  \item Alignment hurts: report it as negative transfer.
\end{ideas}

\section{Role of group members}
\todo{Explicit division of work: who does what.}
\begin{ideas}
  \item Possible work packages: simulation and data; policy and baselines; IDM
    and adaptation; evaluation and integration.
\end{ideas}

{\footnotesize\textbf{Generative-AI disclosure.} Codex (OpenAI) and Claude Code
(Anthropic) assisted with the document structure, starting ideas, literature
search, bibliography, and LaTeX formatting. \todo{Confirm: the proposal text was
written by the team.}}

% Only works cited with \citep{} appear in the references.
\clearpage
\bibliographystyle{unsrtnat}
\bibliography{../references}
\end{document}
